\title{3.5-meter Segmented-Mirror Robotic Space Telescope}
\subtitle{Mission White Paper: Overall Architecture and Scientific Mission}
\newcommand{\wpauthors}{Yong-Woo Kang, Sang Hyun Lee, Juhan Kim, Jeong-Yeol Han, Sungwook E. Hong, Bongkon Moon, Donguk Song, Juhyung Kang, Chung-Uk Lee, Sangmo Tony Sohn, Arman Shafieloo, David Parkinson, Hong Soo Park, Dohyeong Kim, Chan Park, Jungjoo Sohn, Sang Chul Kim, Young-Beom Jeon, Jong-Hak Woo, Myeong-Gu Park}
\date{2026}
%% --- abstract + keywords (use after \maketitle) ---
\newcommand{\wpabstract}{%
A 3.5-meter segmented-mirror robotic space telescope is under study as a space-based observatory for precision astrophysical observations and rapid-response transient astronomy in the 0.2--1.5 μm wavelength range. The telescope adopts a Cassegrain optical configuration optimized to deliver diffraction-limited performance across a wide, flat focal plane, achieving a Strehl ratio greater than 0.8 at 633 nm. The proposed scientific payload includes a Wide-field Camera (WC), a spectroscopic instrument, and an optional Exoplanet Imaging Coronagraph. The Wide-field Camera (WC) provides multi-wavelength imaging and high-cadence time-series photometry over a field of view ranging from 10′×10′ to 30′×30′. The spectroscopic configuration and resolving power remain under study to accommodate the requirements of the principal science programs. An optional Exoplanet Imaging Coronagraph is being investigated for high-contrast imaging of nearby planetary systems, with performance goals extending toward raw contrasts of approximately 10⁻⁸ and improved post-processed performance. Candidate orbital configurations, including Earth orbit and the Sun–Earth L2 region, are currently being evaluated. Planned investigations include gravitational-wave counterparts, rapidly evolving transients, Type Ia supernova cosmology, direct imaging of exoplanets, and exoplanet atmospheric spectroscopy. Although driven by these core scientific objectives, the observatory is conceived as a general-purpose facility providing open-access observing time to the international scientific community. This paper presents the preliminary architecture, performance goals, and scientific mission of the proposed 3.5-meter space telescope.
}
\newcommand{\wpkeywords}{\textbf{Keywords:} segmented mirror, robotic, space telescope, exoplanet, cosmology, time-domain, multi-messenger}
